\documentclass[11pt]{article}
\usepackage[margin=1.2in]{geometry}
\usepackage{amsmath, amssymb}
\usepackage{booktabs}
\usepackage{hyperref}
\usepackage{parskip}
\usepackage{titlesec}
\usepackage{enumitem}

\titleformat{\section}{\large\bfseries}{{\thesection.}}{0.5em}{}
\titleformat{\subsection}{\normalsize\bfseries}{{\thesubsection.}}{0.5em}{}

\title{\textbf{GammaZero: Learned Value Aggregation} \\ \large Design Document}
\author{David Shin}
\date{\today}

\begin{document}

\maketitle
\tableofcontents
\newpage

% ─────────────────────────────────────────────
\section{Motivation}

AlphaZero aggregates child value estimates into a parent estimate using a
visit-count-weighted average:
\[
  Q(p) = \frac{V(p) + \sum_i N_i \, Q_i}{1 + \sum_i N_i}.
\]
This is a reasonable heuristic but has well-known theoretical weaknesses. It
conflates \emph{exploratory} visits (induced by PUCT's exploration term and auxiliary
mechanisms like policy temperature and Dirichlet noise) with
\emph{confirmatory} visits (made because a child is believed to be good).
A child that was heavily explored but ultimately found to be poor will still
exert downward pressure on $Q(p)$ proportional to its visit count. Additionally, a child
suddenly found to be very good needs many more visits to offset highly-visited-but-inferior
siblings than should be theoretically necessary.

GammaZero replaces this hand-crafted aggregation with a \textbf{learned
aggregation network} $\mathcal{A}$ that takes children statistics as input and
outputs a posterior parent estimate $Q(p)$. The hope is that
$\mathcal{A}$ learns to discount exploratory visits and produce better-calibrated
parent estimates, which in turn improves the quality of search.

% ─────────────────────────────────────────────
\section{Architecture Overview}

\subsection{Main Neural Network}

The main NN $f_\theta$ maps a game state $s$ to:
\begin{itemize}[noitemsep]
  \item \textbf{Policy head} $P$: prior move probabilities (same as AlphaZero).
  \item \textbf{Value head} $V$: prior scalar value estimate (same as AlphaZero).
  \item \textbf{Uncertainty head} $U$: prior scalar uncertainty estimate (based
        on KataGo; see Section~\ref{sec:U}).
  \item \textbf{Action-value head} $AV$: per-action value predictions (new; see
        Section~\ref{sec:AV}).
  \item \textbf{Action-uncertainty head} $AU$: per-action uncertainty predictions
        (new; see Section~\ref{sec:AV}).
  \item \textbf{Latent representation} $\mathbf{z}$: a compressed projection of
        the trunk features, stored at the node upon expansion (see
        Section~\ref{sec:z}).
\end{itemize}

The main NN is evaluated \emph{once per node expansion}, exactly as in
AlphaZero.

\subsection{Latent Projection ($\mathbf{z}$)}
\label{sec:z}

The aggregation network conditions on positional context from the main NN. If
the full trunk features were stored at each node, the memory cost would be
prohibitive.
To address this, a small \textbf{projection head} maps the trunk features down
to a compact representation $\mathbf{z}$.

The simplest option is a \textbf{flat} projection to a vector of fixed
dimension $d_z$ (e.g., 32--64), capturing global positional properties such as
tactical sharpness or game phase. An alternative is a \textbf{spatial}
projection that retains board structure; see Section~\ref{sec:spatial-z} for
further discussion.

During training, there is only one network $\mathcal{N}$ comprising the
combination of $f_\theta$ and $\mathcal{A}_\phi$. The portion pertaining to
$\mathcal{A}_\phi$ is merely a head $H_\phi$, emanating from the network
trunk. At export time, $H_\phi$ is detached from the rest of $\mathcal{N}$,
resulting in separate networks $f_\theta$ and $\mathcal{A}_\phi$, with the
$\mathbf{z}$ projection living in $f_\theta$. This means the forward pass of
$f_\theta$ produces the compact $\mathbf{z}$ directly, and only the small
representation is stored at the node. $\mathcal{A}_\phi$ at inference time
never sees the full trunk features.

\subsection{Aggregation Network}
\label{sec:agg-net}

The aggregation network $\mathcal{A}_\phi$ is a lightweight network evaluated
at every \emph{backup step}. Its outputs are stored at the node and read
during selection (see Section~\ref{sec:puct}). The outputs for a node $p$ are:
\begin{itemize}[noitemsep]
  \item $Q(p)$: posterior value estimate for the parent node.
  \item $W(p)$: posterior uncertainty estimate for the parent node.
\end{itemize}

The inputs for a node $p$ with children $\{c_1, \ldots, c_k\}$ are:
\begin{itemize}[noitemsep]
  \item Per-child statistics: $\{Q_i, W_i, N_i\}$ for each child $c_i$.
        For expanded children ($N_i > 0$), $Q_i$ and $W_i$ are outputs of
        $\mathcal{A}_\phi$ previously computed at child $c_i$.
        For unexpanded children ($N_i = 0$), $Q_i = AV_i$ and $W_i = AU_i$,
        the per-action predictions from $f_\theta$ (see Section~\ref{sec:AV}).
  \item The prior outputs of $f_\theta$ at $p$: $\{P(p), V(p), U(p), AV(p), AU(p), 
        \mathbf{z}(p)\}$.
  \item Baseline aggregates derived from the Law of Total Expectation (LoTE) and
        the Law of Total Variance (LoTV):

        \begin{align*}
Q^*(p) &=  \frac{V(p) + \sum_i N_iQ_i}{1 + \sum_i N_i} \\
W^*(p) &= \frac{U(p) + (V(p) - Q^*(p))^2 + \sum_i N_i(W_i + (Q_i - Q^*(p))^2)}{1 + \sum_i N_i}
        \end{align*}

  The LoTV baseline $W^*(p)$ is motivated by interpreting $W$ as a variance.
  This interpretation is loosely justified: the $W$ target is defined as a
  squared difference between $Q$ values (see Section~\ref{sec:U}), which
  matches the structure of variance computation. The interpretation need not be
  rigorous --- $W^*(p)$ serves only as an input suggestion that
  $\mathcal{A}_\phi$ is free to use or ignore.
\end{itemize}

Because $\mathcal{A}_\phi$ is called at every backup step (typically
$O(n \log n)$ times in a search tree with $n$ nodes), it must be small. The
latent $\mathbf{z}$ provides positional context, allowing $\mathcal{A}_\phi$
to modulate its aggregation strategy based on the character of the position
(e.g., tactical sharpness, game phase).

Note that the per-child inputs $Q_i$ and $W_i$ for expanded children are
themselves outputs of $\mathcal{A}_\phi$ from earlier backup steps lower in the
tree. The aggregation network's inputs thus include its own prior outputs,
giving the system a recursive structure. During training, gradients are stopped
at this boundary: the child-level $Q_i$ and $W_i$ values are treated as fixed
inputs.

% ─────────────────────────────────────────────
\section{Training Targets}

\subsection{Target for $Q(p)$}

The training target for $Q(p)$ is the game outcome $Z$. This is the same
supervision signal used for the value head $V$. The framing is:

\begin{itemize}[noitemsep]
  \item $V(p)$ is a \textbf{prior} value estimate, conditioned only on the
        board state.
  \item $Q(p)$ is a \textbf{posterior} value estimate, additionally
        conditioned on observed children statistics from search.
\end{itemize}

Since the network can always represent the function ``output $Q^*(p)$'' as a
special case, it can in principle recover AlphaZero's aggregation. The
hope is that dynamic children statistics carry additional signal about the
true value.

\subsection{Target for $U(p)$}
\label{sec:U}

The main NN includes a dedicated \textbf{uncertainty head} $U(p)$ that
predicts \textbf{uncertainty}: how much the value estimate is
expected to swing between the current position and the end of the game.
For this, we adapt the target that KataGo uses for its Uncertainty-Weighted MCTS
Playouts mechanism\footnote{See: \url{https://github.com/lightvector/KataGo/blob/master/docs/KataGoMethods.md\#uncertainty-weighted-mcts-playouts}}.
Specifically, at turn $T$:
\[
  U(p)_{\text{target}} = \bigl(V(p) - (1-\lambda)\sum_{t > T} Q_{t}\lambda^{t-1-T}\bigr)^2,
\]
where $Q_t$ represents the final value of $Q$ of the tree root at the end of
turn $t$, and where $\lambda$ is a system hyperparameter. Note that the sum ranges over $t > T$, while KataGo's sum ranges
over $t \geq T$. See Section~\ref{sec:Wp} for the rationale.

$U(p)$ is predicted from board features alone, serving as the \emph{prior}
on uncertainty --- exactly analogous to how $V$ is the prior on value.
$U(p)$ is passed as an input to the aggregation network, providing a baseline
that the posterior $W(p)$ can refine.

The critical grounding property of this target is that $U(p) = 0$ for any
position whose game-theoretic value is fully resolved --- either because it
directly precedes a terminal state, or because its subtree has been
exhaustively searched (e.g., proven mate, proven draw). This gives the
aggregation network a reliable anchor: $W_i = 0$ means child $c_i$'s value
$Q_i$ is completely trustworthy, and $W_i = \epsilon$ means it is nearly so.

As noted in Section~\ref{sec:agg-net}, we interpret $U(p)$ as a variance
of $V(p)$ to compute a LoTV aggregation over children statistics. This is
merely a chosen \textit{interpretation} and is not rigorously justified.
It need not be rigorously justified, as the LoTV-computed aggregation serves
merely as a \textit{suggestion} to $\mathcal{A}_\phi$ on how to compute the
actual uncertainty-aggregation, which $\mathcal{A}_\phi$ is free to use or
ignore. Whether the target definition or the LoTV aggregation are useful
are empirical questions; what ultimately matters is
whether the choices help $\mathcal{A}_\phi$ produce better $Q(p)$ estimates.

\subsection{Target for $W(p)$}
\label{sec:Wp}

$W(p)$ is the posterior version of $U(p)$, output by the aggregation network.
It is trained to predict an analogous target to $U(p)$, but with the benefit of
access to children statistics. Denoting by $Q^{(n)}(p)$ the output of
$\mathcal{A}_\phi$ after $n$ search iterations:
\[
  W(p)_{\text{target}} = \bigl(Q^{(n)}(p) - (1-\lambda)\sum_{t > T} Q_{t}\lambda^{t-1-T}\bigr)^2,
\]

As for why our sum ranges over $t > T$ rather than $t \geq T$: the expected
difference between the current $Q^{(n)}(p)$ and $Q_T$ will depend on $n$, the
search iteration for which the sample was obtained. This would introduce an
inappropriate coupling between $n$ and $W(p)_\text{target}$.

The LoTV-derived baseline $W^*(p)$ serves as a
regularization anchor during early training. See Section~\ref{sec:bootstrap}.

\subsection{Targets for $AV$ and $AU$}
\label{sec:AV}

The main NN includes per-action heads $AV(s, a)$ and $AU(s, a)$ that predict
the value and uncertainty at the child state reached by action $a$. Their
training targets are the main NN's own evaluations at those child states:
\begin{align*}
  AV(s, a)_{\text{target}} &= V(s'), \\
  AU(s, a)_{\text{target}} &= U(s'),
\end{align*}
where $s'$ is the state reached by playing action $a$ from $s$.

To produce these targets, children of sampled root nodes are expanded during
training. This can be done exhaustively (expanding all children) or by sampling
a random subset and computing loss only on that subset.

$AV$ and $AU$ serve as the default $Q_i$ and $W_i$ inputs for unexpanded
children when feeding into $\mathcal{A}_\phi$. This allows the aggregation
network to reason about all children --- not just those that have been visited
--- using the main NN's prior predictions as stand-ins.

Note that AlphaZeroArcade\footnote{\url{https://github.com/shindavid/AlphaZeroArcade}}
already uses an $AV$ head, whose output is used for the $Q$ term in the PUCT
formula for unvisited children.

\subsection{Prior-Consistency Regularization}
\label{sec:prior_consistency}

When all children are unexpanded ($N_i = 0$ for all $i$), the aggregation
network has no search evidence to incorporate. The only child-level information
available is the $AV$/$AU$ priors. In this case, $\mathcal{A}_\phi$ has no
information beyond what $f_\theta$ already captured, so the posterior should
equal the prior:
\[
  Q(p) = V(p), \qquad W(p) = U(p) \qquad \text{when all } N_i = 0.
\]
This is enforced as a regularization term during training: a fraction of
training samples are presented with all $N_i = 0$, and a loss term encourages
$\mathcal{A}_\phi$ to reproduce $V(p)$ and $U(p)$ in this regime.

% ─────────────────────────────────────────────
\section{Training Procedure}

\subsection{Sampling for $\mathcal{A}_\phi$}

Following KataGo's Playout Cap Randomization, some fraction of root nodes
are sampled for NN training. Of these, a configurable fraction $\alpha$ is
allocated to training the main NN $f_\theta$, and the remainder $1 - \alpha$
to training the aggregation NN $\mathcal{A}_\phi$. The split ratio $\alpha$
is a system hyperparameter; since $f_\theta$ is a larger network learning a
harder function, $\alpha > 0.5$ is likely appropriate.

For positions allocated to $\mathcal{A}_\phi$ training, if $N$ is the search
budget in iterations, we uniformly sample some $n$ from $\{1, \ldots, N\}$
and snapshot the tree statistics after $n$ iterations have been performed.
This means $\mathcal{A}_\phi$ learns to produce a useful posterior estimate
given any amount of search evidence, from very little ($n$ small) to nearly
the full budget ($n \approx N$).

\subsection{Bootstrapping and Regularization}
\label{sec:bootstrap}

Early in training, $\mathcal{A}_\phi$ is uninformed, which risks destabilizing
the feedback loop between search quality and training targets. To mitigate
this, the network is regularized during early training to match the
hand-crafted baselines:

\begin{align*}
  Q(p)_{\text{baseline}} &= Q^*(p), \\
  W(p)_{\text{baseline}} &= W^*(p).
\end{align*}

This regularization is annealed away as training progresses. Because $Q^*(p)$
and $W^*(p)$ are inputs into $\mathcal{A}_\phi$, they are representable and
thus constitute reachable regularization targets. The network will learn
initially to simply directly output these inputs, providing some initial
stability. As the annealing tapers, the network should gradually mature to
produce more nuanced outputs.

The prior-consistency regularization (Section~\ref{sec:prior_consistency})
provides an additional stability mechanism: even if the LoTE/LoTV
regularization is annealed away, $\mathcal{A}_\phi$ retains sensible behavior
in the zero-search regime.

% ─────────────────────────────────────────────
\section{Inference and Search Integration}

\subsection{Computational Cost}

The main NN is evaluated once per node expansion, unchanged from AlphaZero.
The aggregation network $\mathcal{A}_\phi$ is evaluated at every backup step.
Its $Q(p)$ and $W(p)$ outputs are stored at the node and read during
selection, requiring no additional evaluation at selection time. Inference
cost is kept low by:
\begin{itemize}[noitemsep]
  \item Keeping $\mathcal{A}_\phi$ small (few linear layers).
  \item Reusing the stored latent $\mathbf{z}(p)$ rather than re-running the
        main NN.
\end{itemize}

In self-play mode, evaluations of $\mathcal{A}_\phi$ across simultaneous games
are batched, amortizing cost. In single-game mode, we may be able to adopt a
similar mechanism to standard MCTS parallelization, which involves multiple
search threads and virtual visits. Some details remain to be worked out.

It may be possible to incorporate
NNUE\footnote{See \url{https://www.chessprogramming.org/NNUE}}
mechanics. In a game like chess, the policy space is large, while only a sparse
subset corresponds to legal moves in any given position. Furthermore, typically
only one child will have changes to its statistics during aggregation (more than
one is possible due to multithreaded search or due to move transpositions).
These properties are the ingredients that make NNUE feasible in chess and
Othello.

\subsection{Use in PUCT}
\label{sec:puct}

$\mathcal{A}_\phi$'s stored $Q_i$ output for each child $c_i$ is used
directly in the PUCT selection formula:
\[
  \text{PUCT}_i = Q_i + c \cdot P_i \cdot \frac{\sqrt{1 + N(p)}}{1 + N_i}.
\]
This is a departure from standard AlphaZero, where $Q_i$ during selection is
simply the running average $\text{total\_value}_i / N_i$. Here, $Q_i$ is the
aggregation network's learned posterior estimate, computed during backup and
cached at the node.

\subsection{Future Direction: Replacing PUCT}

Later, we may add a head to $\mathcal{A}_\phi$ that outputs a score that
replaces PUCT entirely. The challenge is devising a training target that
captures the purpose of PUCT.

One idea for this is to define a notion of ``exploration worthiness''.
PUCT aims to identify actions worth exploring. What makes an action
worth exploring? One candidate definition is to say that an action is
worth exploring if doing so has strong potential for changing our current $Q$
estimate. This suggests using a training target that is based on the
change in $Q(p)$ induced by visiting a given child. A similar alternative 
is to say that an action is worth exploring if \textit{not} doing so for
the remaining search budget is likely to produce a different $Q$
estimate than an unconstrained usage of the remaining search budget.
Still another possibility is to locate exploration worthiness in a likely
reduction of $W(p)$, rather than an expected delta in $Q(p)$. 
It may be possible to produce targets corresponding
to these notions.

All these options have the property that we can use an exact score of 0 for
fully resolved positions (e.g., proven mate, proven draw), while a NN estimated
position can be made to return a strictly positive score. This has the nice
property that the selection criterion would naturally avoid re-exploring 
fully resolved positions, a property that standard PUCT lacks.

If something can be found to work here, it may be possible to remove one or more
heads among $P$, $AU$, and $AV$, as they exist primarily to facilitate the
PUCT criterion. 

% ─────────────────────────────────────────────
\section{Spatial Latent Representations}
\label{sec:spatial-z}

Section~\ref{sec:z} introduced the latent projection $\mathbf{z}$ as a
compressed representation of the trunk features. The simplest design is a flat
vector of dimension $d_z$, which captures global positional properties but
discards spatial structure. An alternative is to retain the board geometry by
projecting to a \textbf{spatial tensor} of shape $(H, W, r)$, where $H \times
W$ matches the board dimensions and $r$ is a small channel depth (e.g., 4--8).

A spatial $\mathbf{z}$ enables $\mathcal{A}_\phi$ to represent
\textbf{correlation beliefs between actions}. For example, in Go on a $19
\times 19$ board, a spatial projection of shape $(19, 19, r)$ could learn to
map each candidate move to a vector on an $r$-dimensional hypersphere, with the
cosine similarity between two move-vectors encoding the aggregation network's
belief about how correlated their values are. This would allow
$\mathcal{A}_\phi$ to reason about questions like: ``move $A$ and move $B$
both attack the same group, so if $A$'s value dropped after search, $B$'s
value likely dropped too.'' A flat $\mathbf{z}$ cannot represent such
per-action structure.

The memory cost of a spatial $\mathbf{z}$ is $(H \times W \times r)$ floats
per node. For Go with $r = 8$, this is $19 \times 19 \times 8 = 2{,}888$
floats --- reasonable for practical use.

This design is more natural for games with spatially structured action spaces
(Go, Hex, Othello) than for games like Chess, where the action space is a set
of (source, destination) pairs rather than single board intersections.

Whether spatial $\mathbf{z}$ actually helps is an empirical question. Even if
the representation has the capacity to encode correlations, the aggregation
network must learn to use them, and the training signal (game outcomes) may be
too diffuse to drive fine-grained correlation learning. The flat design is the
safe default; spatial $\mathbf{z}$ is an avenue for games where it is natural.

% ─────────────────────────────────────────────
\section{Stochastic Games}

For games with explicit stochastic events whose probability distributions are
known (e.g., dice rolls), the uncertainty target conflates epistemic uncertainty
with aleatory uncertainty from the game mechanics. To compensate, a
\textbf{luck correction} is applied using the main NN's value predictions at
chance nodes.

In games with stochastic events, the search tree contains \textbf{chance
nodes} (representing stochastic outcomes) alongside \textbf{player nodes}
(representing player decisions). The main NN evaluates both node types,
producing $V$ predictions at each. At a chance node with outcome $\omega$,
the luck delta is:
\[
  \delta_\omega = V(\text{child}) - V(\text{parent}),
\]
i.e., the shift in value prediction between the pre-event state and the
post-event state. If per-action value heads are available, this simplifies to
$\delta_\omega = V(\text{child}) - AV(\text{parent}, a_\omega)$, where
$a_\omega$ is the action corresponding to outcome $\omega$.

The luck-adjusted $Q$ along a trajectory is then:
\[
  Q_{\text{adj}} = Q_{\text{raw}} - \sum_{\text{events}} \delta_\omega.
\]
This is a martingale correction: it removes the random-walk component from the
value trajectory, leaving only the skill/position component. The $U(p)$ and
$W(p)$ training targets are then computed against $Q_{\text{adj}}$ rather than
$Q_{\text{raw}}$.

No separate luck estimator network is required --- the correction falls out of
the main NN's existing value predictions at chance nodes.

For deterministic games (Chess, Go, Hex, Connect4), there are no chance nodes
and this section does not apply.

% ─────────────────────────────────────────────
\section{Hyperparameters}

Table~\ref{tab:hyperparams} summarizes the tunable hyperparameters of the
system.

\begin{table}[h]
\centering
\begin{tabular}{@{}lll@{}}
\toprule
\textbf{Symbol} & \textbf{Description} & \textbf{Default / Suggested} \\
\midrule
$\alpha$ & Fraction of sampled roots allocated to $f_\theta$ training
         & $> 0.5$ (TBD) \\
$\lambda$ & Exponential decay for uncertainty targets
          & $5/6$ (from KataGo) \\
$d_z$    & Projection dimension (flat $\mathbf{z}$)
         & 32--64 \\
$r$      & Channel depth (spatial $\mathbf{z}$)
         & 4--8 \\
         & Annealing schedule for LoTE/LoTV regularization
         & TBD \\
         & Fraction of $N_i = 0$ samples for prior-consistency
         & TBD \\
\bottomrule
\end{tabular}
\caption{System hyperparameters.}
\label{tab:hyperparams}
\end{table}

% ─────────────────────────────────────────────
\section{Open Questions}

\begin{enumerate}
  \item \textbf{Aggregation network architecture and inference mechanics.}
        How many layers? NNUE? How does query batching work in single-game
        mode?
  \item \textbf{Regularization details.} There are multiple components of
        annealing and regularization (LoTE/LoTV baseline, prior-consistency,
        and potentially independent annealing rates for $Q$ vs.\ $W$). The
        interactions between these components need to be explored empirically.
  \item \textbf{$AV$/$AU$ training efficiency.} Expanding all children at
        sampled roots may be expensive for games with large branching factors.
        The subset-sampling alternative introduces variance; the right balance
        is to be determined.
  \item \textbf{Empirical validation} on Connect4, then Othello, Hex, Chess.
\end{enumerate}

\end{document}
